% Test of code side by side (didactic): \codebycode, \runbyrun and
% didactic_pair, in the notes (memoir + beamerarticle).
%
% Build: see the Makefile (pdflatex -shell-escape, pythontex, pdflatex twice).
%
% Expected in the resulting PDF:
%   1  The footnote before the first pair (BEFORENOTE) is printed, and so are
%      BETWEENNOTE and AFTERNOTE.
%   2  The two main functions side by side over the full width, the margin
%      included, with pairgreet1.py and pairgreet2.py above the columns.
%      pairgreet1.py calls main() at top level and asks for input: it must
%      not be run (the build doesn't stop for input).
%   3  greet without its docstring beside Person.greet without its
%      docstring, the right column labelled Person.greet.
%   4  The whole phonebook dictionary (from/to with braces) beside lines
%      4-8 of pairgreet2.py (the class Person up to its __init__).
%   5  \runbyrun: "Hello, world!" beside the shell's "hej".
%   6  Three columns from didactic_pair, labelled one, two and three.
\documentclass[a4paper]{memoir}
\usepackage{beamerarticle}
\usepackage{lipsum}
\usepackage{didactic}
\begin{document}
\section{Pairs}
\lipsum[1][1-3] Before the pair\footnote{BEFORENOTE must survive.} more text.

\codebycode[name=main,label]{python}{pairgreet1.py}%
  [name=main,label]{python}{pairgreet2.py}

Between\footnote{BETWEENNOTE here.} the pairs.

\codebycode[name=greet,nodocstring]{python}{pairgreet1.py}%
  [name=Person.greet,nodocstring,label={\texttt{Person.greet}}]%
  {python}{pairgreet2.py}

\codebycode[from={phonebook = \{},to=\}]{python}{pairgreet1.py}%
  [lines=4-8]{python}{pairgreet2.py}

\runbyrun{hello.py}[shell]{echo hej}

\begin{pycode}
didactic_pair("a = 1", "b = 2\nc = 3", "d = 4",
              labels=["one", "two", "three"])
\end{pycode}
After\footnote{AFTERNOTE here.} the end.
\end{document}
